\section{Conclusions and Future Work}
\label{sec:conclusions}

\subsection{Summary of Results}

We have presented a self-contained derivation of the gas-driven permeation
timelag for a bilayer composite membrane using Laplace-transform methods.
The principal results are:

\begin{enumerate}
  \item \textbf{Composite timelag formula.}
    The downstream timelag for a bilayer membrane with layer thicknesses
    $\Ls{1}$, $\Ls{2}$, diffusivities $\Ds{1}$, $\Ds{2}$, and solubilities
    $\Ss{1}$, $\Ss{2}$ is given by \cref{eq:tlag-composite}.  This reduces to
    the classical $L^2/(6D)$ result for a homogeneous membrane and agrees with
    the dimensionless form of \textcite{TranMeldonSontag2021}.

  \item \textbf{Direction independence of the downstream timelag.}
    The downstream timelag is independent of the direction of permeation
    (\cref{prop:direction-independence}).  This has the important practical
    consequence that measuring the downstream timelag in both the forward and
    reverse directions does not yield two independent equations for
    parameter extraction.

  \item \textbf{Direction dependence of the upstream lead time.}
    In contrast to the downstream timelag, the upstream lead time
    $\theta_{\mathrm{up}}$ is direction-dependent
    (\cref{prop:upstream-direction}), providing additional independent
    information.

  \item \textbf{Direction independence of the steady-state flux.}
    Like the timelag, the steady-state flux through a bilayer is independent
    of the permeation direction, depending only on the layer permeabilities
    $\Phi_i = \Ds{i}\Ss{i}$ and $\sqrt{P}$.

  \item \textbf{Extraction methods.}
    We showed that a single permeation experiment on a bilayer yields at most
    three independent quantities ($\tlag$, $\Jss$, $\theta_{\mathrm{up}}$),
    which are insufficient to determine all four unknowns ($\Ds{1}$, $\Ds{2}$,
    $\Ss{1}$, $\Ss{2}$) without additional information.  Three strategies are
    proposed (\cref{sec:extraction}):
    \begin{itemize}
      \item \emph{Method~1}: Single-layer calibration of one constituent
            material, combined with a bilayer experiment.
      \item \emph{Method~2}: Two bilayer samples with different thickness
            ratios, requiring no single-layer data.
      \item \emph{Method~3}: Full transient curve fitting to extract all
            parameters from a single experiment.
    \end{itemize}
    When solubilities are known a priori, the problem reduces to two unknowns
    and a single experiment suffices.
    The extraction is well-conditioned when both layers contribute comparably
    to the total permeation resistance (\cref{prop:well-posed}).
\end{enumerate}

\subsection{Limitations}

The analysis presented here rests on several assumptions:
\begin{itemize}
  \item Fickian diffusion with constant, concentration-independent diffusivities
        in each layer.
  \item Ideal boundary conditions: instantaneous equilibration at the external
        surfaces (no surface resistance).
  \item Ideal interface: thermodynamic equilibrium and flux continuity at the
        layer interface (no interfacial resistance).
  \item Known layer thicknesses $\Ls{1}$, $\Ls{2}$ (solubilities need not be
        known a priori, but additional measurements are then required).
\end{itemize}
Departures from these assumptions---such as concentration-dependent
diffusivity, surface kinetic barriers, interfacial resistance, or hydrogen
trapping---will modify the timelag expressions and may require extended
models.

\subsection{Future Directions}

Several extensions of the present work are of practical interest:
\begin{enumerate}
  \item \textbf{Three or more layers.}
    The Laplace-transform approach generalises naturally to $N$-layer
    composites via transfer-matrix methods, yielding timelag expressions
    as functions of $2N$ material parameters ($\Ds{i}$, $\Ss{i}$).
  \item \textbf{Cylindrical and spherical geometries.}
    Pipeline and pressure vessel applications involve cylindrical membranes,
    where the curvature modifies the steady-state concentration profile and
    the timelag expression.
  \item \textbf{Concentration-dependent diffusivity.}
    In many polymers and some metals, $D$ depends on concentration.  The
    timelag then depends on the upstream concentration $C_0$, and
    multiple experiments at different pressures can provide additional
    constraints.
  \item \textbf{Trapping effects.}
    In metals, hydrogen trapping at microstructural defects increases the
    effective timelag.  Incorporating McNabb--Foster or Oriani trapping models
    into the bilayer framework is an important extension.
  \item \textbf{Numerical validation.}
    The analytical results should be validated against finite-element
    simulations (e.g., using FESTIM or COMSOL) with realistic material
    parameters, providing confidence in both the derivations and the
    extraction algorithms under practical noise levels.
\end{enumerate}
